\subsection{比较尺度}

设 E 是一个被具有基 F 的滤子所滤过的集，K 是一个非离散赋值域（最常见的是 K = R 或 K = C）。在 \(\mathcal{H}(\mathfrak{F}, \mathrm{K})\) 中模 \(R_{\infty}\) 不等价于 0 的函数（即在 F 的每个集中至少有一点使函数不为零的那些函数）所成的集中，关系 “ \(f \ll g\) 或 f = g” 是一个序关系。

定义 1. 我们说 \(\mathcal{H}(\mathfrak{F},\mathrm{K})\) 的一个由模 \(R_{\infty}\) 不等价于 0 的函数组成的子集 E 是一个比较尺度，如果 E 被关系 “ \(f \ll g\) 或 f = g” 全序化。

换句话说，若 \(f\) 与 \(g\) 是 \(\mathcal{E}\) 中的函数，则关系 \(f \ll g\) 、 \(g \ll f\) 、 \(f = g\) 中总有一个（且只有一个）成立。由此推出，在 \(\mathcal{E}\) 上关系 \(f \asymp g\) （更不用说 \(|f| \sim a |g|\) ，其中 \(a\) 是一个数 \(> 0\) ）蕴含 \(f = g\) 。

一个比较尺度的每个子集显然也是一个比较尺度。

例. 1) 对实的且趋于 \(+\infty\) 的 x，函数 \(x^{\alpha}\) （ \(\alpha\) 为任意实数）所成的集是一个比较尺度。当 F 是以 a 为左端点的开区间所成的集时，函数 \((x - a)^{\alpha}\) 亦然。

\begin{enumerate}
\def\labelenumi{\arabic{enumi})}
\setcounter{enumi}{1}
\item
  当复数 \(z\) 趋于 \(\infty\) 时，函数 \(z^n\) （ \(n\) 为有理整数）所成的集是一个比较尺度；函数 \((z - a)^n\) 亦然（当 \(\mathfrak{F}\) 是该点的邻域滤子在点 \(a \in \mathbf{C}\) 的余集上的迹时）。
\item
  设 F 是赋范空间；当 F 是该点的邻域滤子在 a 的余集上的迹时，函数族 \(\|\mathbf{x}-\mathbf{a}\|^{\alpha}\) （ \(\alpha\) 为任意实数）是一个比较尺度。注意，若 p 与 q 是 F 上两个不同的范数，则两个比较尺度 \(\left(p(\mathbf{x}-\mathbf{a})\right)^{\alpha}\) 与 \(\left(q(\mathbf{x}-\mathbf{a})\right)^{\alpha}\) 的并一般不是一个比较尺度。
\item
  当实数 \(x\) 趋于 \(+\infty\) 时，函数族 \(\mathcal{E}\) 由形如 \(\exp(p(x))\) 的函数（其中 \(p\) 取遍无常数项的（实系数）多项式所成的集）组成，它是一个比较尺度：只需注意， \(\mathcal{E}\) 中两个函数的商仍在 \(\mathcal{E}\) 中，且函数 \(\exp(p(x))\) 必趋于 0 或 \(+\infty\) （如果 \(p \neq 0\) ）；事实上， \(p(x) \sim \alpha x^n\) （其中 \(n > 0\) 且 \(\alpha \neq 0\) ）；若 \(\alpha > 0\) ，则有 \(p(x) > \frac{1}{2}\alpha x^n\) （对充分大的 \(x\) ）；若 \(\alpha < 0\) ，则有 \(p(x) < \frac{1}{2}\alpha x^n\) （对充分大的 \(x\) ）；在第一种情形 \(\exp(p(x))\) 趋于 \(+\infty\) ，在第二种情形趋于 0。
\item
  对实的且趋于 \(+\infty\) 的 x，形如 \(x^{\alpha}(\log x)^{\beta}\) （对 x \textgreater{} 1 有定义）的函数所成的集 E（其中 \(\alpha\) 与 \(\beta\) 为任意实数）是一个比较尺度。事实上，这里同样，E 中两个函数的商是 E 中的一个函数；只需证明：若 \(\alpha\) 与 \(\beta\) 不全为零，则 \(x^{\alpha}(\log x)^{\beta}\) 趋于 0 或 \(+\infty\) ；这在 \(\alpha = 0, \beta \neq 0\) 时是明显的；若 \(\alpha > 0\) ，有 \((\log x)^{-\beta} \ll x^{\alpha}\) ，而若 \(\alpha < 0\) ，有 \((\log x)^{\beta} \ll x^{-\alpha}\) （对任何 \(\beta\) ），由此得命题。
\end{enumerate}

注意，最后这个比较尺度是一个全序集（对关系 “ \(f \ll g\) 或 \(f = g\) ”），其序结构同构于 \(R^{2}\) 上的字典序 (Set Theory, III, p.~157)；回忆在这个结构中关系 \((\alpha, \beta) < (\gamma, \delta)\) 表示 “ \(\alpha < \gamma\) ，或 \(\alpha = \gamma\) 且 \(\beta < \delta\) ”。

类似地，由函数 \(\exp(p(x))\) （其中 p 取遍无常数项多项式所成的集 \(P_{0}\) ）组成的尺度具有同构于 \(P_{0}\) 的序结构的序结构，在其中关系 p \textless{} q 表示多项式 q - p 的主导项具有 \textgreater{} 0 的系数 (cf.~Alg., VI. 19, 例 2)。

设 \(\varphi\) 是从集 \(F\) 到 \(E\) 中的一个映射，使得 \(\overset{-1}{\varphi}(\mathfrak{F})\) 是 \(F\) 上的一个滤子基。若 \(\mathcal{E}\) 是 \(E\) 上的一个比较尺度（对滤子基 \(\mathfrak{F}\) ），则函数 \(f \circ \varphi\) （当 \(f\) 取遍 \(\mathcal{E}\) ）在 \(F\) 上组成一个比较尺度（对滤子基 \(\overset{-1}{\varphi}(\mathfrak{F})\) ）。

